\documentclass[12pt,a4paper]{article}

% ---- Paquetes ----
\usepackage[utf8]{inputenc}
\usepackage[T1]{fontenc}
\usepackage[spanish]{babel}
\usepackage{geometry}
\usepackage{hyperref}
\usepackage{xcolor}
\usepackage{listings}
\usepackage{enumitem}
\usepackage{float}
\usepackage{booktabs}
\usepackage{longtable}
\usepackage{caption}

\geometry{margin=2.5cm}

% ---- Configuración de listings ----
\lstset{
    language=Python,
    basicstyle=\small\ttfamily,
    keywordstyle=\color{blue}\bfseries,
    stringstyle=\color{red},
    commentstyle=\color{green!50!black},
    frame=single,
    numbers=left,
    numberstyle=\tiny,
    breaklines=true,
    showstringspaces=false,
    tabsize=4,
    captionpos=b,
    backgroundcolor=\color{gray!5},
    literate=
        {á}{{\'a}}1 {é}{{\'e}}1 {í}{{\'i}}1 {ó}{{\'o}}1 {ú}{{\'u}}1
        {ñ}{{\~n}}1 {ü}{{\"u}}1 {Á}{{\'A}}1 {É}{{\'E}}1 {Í}{{\'I}}1
        {Ó}{{\'O}}1 {Ú}{{\'U}}1
}

% Estilo para JSON
\lstdefinelanguage{json}{
    basicstyle=\small\ttfamily,
    stringstyle=\color{teal},
    keywordstyle=\color{blue},
    morestring=[b]",
    literate=
        {á}{{\'a}}1 {é}{{\'e}}1 {í}{{\'i}}1 {ó}{{\'o}}1 {ú}{{\'u}}1
        {ñ}{{\~n}}1
}

\lstdefinestyle{bash}{
    language=bash,
    basicstyle=\small\ttfamily\color{black},
    backgroundcolor=\color{black!5},
    frame=single,
    numbers=none,
    breaklines=true,
}

% ---- Metadata ----
\title{\textbf{Endpoints REST: Actualización y Eliminación}\\de Obras de Arte\\\large Proyecto Atrium -- Museo de Arte Contemporáneo}
\author{Documentación Técnica}
\date{\today}

\begin{document}

\maketitle
\thispagestyle{empty}
\newpage

\tableofcontents
\newpage

% ===================================================================
\section{Introducción}
% ===================================================================

El sistema \textbf{Atrium} expone dos operaciones fundamentales para la
gestión de obras de arte: actualización (\texttt{PUT}) y eliminación
(\texttt{DELETE}). Estos endpoints permiten modificar los datos de una obra
existente o removerla del catálogo, respectivamente.

El presente documento detalla el funcionamiento interno de ambos endpoints,
incluyendo su implementación en la capa de rutas y controladores, la
validación de datos mediante Pydantic, los códigos de respuesta, y cómo
realizar pruebas con \texttt{curl}.

\subsection{Tecnologías utilizadas}

\begin{itemize}
    \item \textbf{FastAPI} (0.136.1) -- Framework web asíncrono.
    \item \textbf{Pydantic v2} (2.13.4) -- Validación y serialización de datos.
    \item \textbf{MongoDB} con \textbf{Motor} -- Base de datos NoSQL.
    \item \textbf{BSON} (\texttt{objectid}) -- Identificadores de documentos.
    \item \textbf{Python 3.12+} con tipado estático.
\end{itemize}

\subsection{Estructura de archivos involucrados}

\begin{verbatim}
app/
  routes/
    artworks.py           # Definicion de los endpoints
  controllers/
    artwork_controller.py # Logica de negocio (CRUD contra MongoDB)
  schemas/
    artwork.py            # Modelos Pydantic y validacion
\end{verbatim}

% ===================================================================
\section{Endpoint PUT: Actualizar obra}
% ===================================================================

\subsection{Definición}

\begin{lstlisting}[language=json]
PUT /artworks/{artwork_id}
Content-Type: application/json
\end{lstlisting}

\subsection{Parámetros}

\begin{longtable}{lll}
\toprule
\textbf{Parámetro} & \textbf{Tipo} & \textbf{Descripción} \\
\midrule
\endfirsthead
\toprule
\textbf{Parámetro} & \textbf{Tipo} & \textbf{Descripción} \\
\midrule
\endhead
\bottomrule
\endfoot
\texttt{artwork\_id} (path) & \texttt{str} & ObjectId de MongoDB (24 caracteres hex) \\
\texttt{Body} & \texttt{ArtworkCreate} & JSON con los campos a actualizar \\
\bottomrule
\end{longtable}

\subsection{Comportamiento}

El endpoint recibe un identificador de obra y un cuerpo JSON con los campos
a actualizar. El controlador:

\begin{enumerate}
    \item Convierte el \texttt{artwork\_id} a un \texttt{ObjectId} de BSON.
    \item Si el ID no es válido (formato incorrecto), retorna \texttt{404}.
    \item Elimina el campo \texttt{\_id} del cuerpo por seguridad.
    \item Convierte los campos de tipo \texttt{date} a \texttt{datetime} para MongoDB.
    \item Asigna la marca de tiempo \texttt{updated\_at} con la hora UTC actual.
    \item Ejecuta \texttt{find\_one\_and\_update} con \texttt{\$set}.
    \item Si el documento no existe, retorna \texttt{404}.
    \item Si existe, retorna el documento actualizado.
\end{enumerate}

\subsection{Código de la ruta}

\begin{lstlisting}
@router.put("/{artwork_id}")
async def update_artwork(artwork_id: str, payload: ArtworkCreate = Body(...)):
    data = payload.model_dump()
    result = await update_artwork_db(artwork_id, data)
    if result is None:
        raise HTTPException(
            status_code=status.HTTP_404_NOT_FOUND,
            detail="Obra no encontrada",
        )
    result["_id"] = str(result["_id"])
    return TypeAdapter(ArtworkResponse).validate_python(result)
\end{lstlisting}

\subsection{Código del controlador}

\begin{lstlisting}
async def update_artwork_db(artwork_id: str, data: dict[str, Any]) -> dict[str, Any] | None:
    try:
        obj_id = ObjectId(artwork_id)
    except errors.InvalidId:
        return None

    data.pop("_id", None)
    _prepare_for_mongo(data)
    data["updated_at"] = datetime.utcnow()

    result = await db[COLLECTION].find_one_and_update(
        {"_id": obj_id},
        {"$set": data},
        return_document=True,
    )
    return result
\end{lstlisting}

\subsection{Códigos de respuesta}

\begin{longtable}{lll}
\toprule
\textbf{Código} & \textbf{Descripción} & \textbf{Cuerpo} \\
\midrule
\endfirsthead
\toprule
\textbf{Código} & \textbf{Descripción} & \textbf{Cuerpo} \\
\midrule
\endhead
\bottomrule
\endfoot
\texttt{200 OK} & Obra actualizada exitosamente & \texttt{ArtworkResponse} (JSON) \\
\texttt{404 Not Found} & ID inválido o obra inexistente & \texttt{\{"detail": "Obra no encontrada"\}} \\
\texttt{422 Unprocessable Entity} & Body no cumple la validación & Error de validación Pydantic \\
\bottomrule
\end{longtable}

% ===================================================================
\section{Endpoint DELETE: Eliminar obra}
% ===================================================================

\subsection{Definición}

\begin{lstlisting}[language=json]
DELETE /artworks/{artwork_id}
\end{lstlisting}

\subsection{Parámetros}

\begin{longtable}{lll}
\toprule
\textbf{Parámetro} & \textbf{Tipo} & \textbf{Descripción} \\
\midrule
\endfirsthead
\toprule
\textbf{Parámetro} & \textbf{Tipo} & \textbf{Descripción} \\
\midrule
\endhead
\bottomrule
\endfoot
\texttt{artwork\_id} (path) & \texttt{str} & ObjectId de MongoDB (24 caracteres hex) \\
\bottomrule
\end{longtable}

\subsection{Comportamiento}

\begin{enumerate}
    \item Convierte el \texttt{artwork\_id} a \texttt{ObjectId}.
    \item Si el ID no es válido, retorna \texttt{404}.
    \item Ejecuta \texttt{delete\_one} en la colección.
    \item Si ningún documento fue eliminado (\texttt{deleted\_count == 0}), retorna \texttt{404}.
    \item Si se eliminó correctamente, retorna \texttt{204 No Content}.
\end{enumerate}

\subsection{Código de la ruta}

\begin{lstlisting}
@router.delete("/{artwork_id}", status_code=status.HTTP_204_NO_CONTENT)
async def delete_artwork(artwork_id: str):
    deleted = await delete_artwork_db(artwork_id)
    if not deleted:
        raise HTTPException(
            status_code=status.HTTP_404_NOT_FOUND,
            detail="Obra no encontrada",
        )
\end{lstlisting}

\subsection{Código del controlador}

\begin{lstlisting}
async def delete_artwork_db(artwork_id: str) -> bool:
    try:
        obj_id = ObjectId(artwork_id)
    except errors.InvalidId:
        return False

    result = await db[COLLECTION].delete_one({"_id": obj_id})
    return result.deleted_count > 0
\end{lstlisting}

\subsection{Códigos de respuesta}

\begin{longtable}{lll}
\toprule
\textbf{Código} & \textbf{Descripción} & \textbf{Cuerpo} \\
\midrule
\endfirsthead
\toprule
\textbf{Código} & \textbf{Descripción} & \textbf{Cuerpo} \\
\midrule
\endhead
\bottomrule
\endfoot
\texttt{204 No Content} & Obra eliminada exitosamente & \emph{(vacío)} \\
\texttt{404 Not Found} & ID inválido o obra inexistente & \texttt{\{"detail": "Obra no encontrada"\}} \\
\bottomrule
\end{longtable}

% ===================================================================
\section{Validación de datos (ArtworkCreate)}
% ===================================================================

Ambos endpoints utilizan el mismo esquema de validación \texttt{ArtworkCreate},
una unión discriminada de cinco modelos según el campo \texttt{genero}.

\subsection{Campos comunes (ArtworkBase)}

\begin{longtable}{lll}
\toprule
\textbf{Campo} & \textbf{Tipo} & \textbf{Validación} \\
\midrule
\endfirsthead
\toprule
\textbf{Campo} & \textbf{Tipo} & \textbf{Validación} \\
\midrule
\endhead
\bottomrule
\endfoot
\texttt{nombre} & \texttt{str} & 1--200 caracteres, obligatorio \\
\texttt{artista\_id} & \texttt{str} & Referencia al artista, obligatorio \\
\texttt{precio\_venta} & \texttt{float} & $> 0$, redondeado a 2 decimales \\
\texttt{fecha\_creacion} & \texttt{date} & No puede ser futura \\
\texttt{foto} & \texttt{str} & URL de la imagen, obligatorio \\
\texttt{descripcion} & \texttt{Optional[str]} & Máx. 2000 caracteres, opcional \\
\bottomrule
\end{longtable}

\subsection{Campos específicos por género}

\subsubsection{Pintura}

\begin{lstlisting}[language=json]
{
    "genero": "pintura",
    "tecnica": "oleo | acrilico | acuarela | tempera",
    "soporte": "lienzo | madera | papel | carton",
    "alto_cm": 73.7,
    "ancho_cm": 92.1
}
\end{lstlisting}

\subsubsection{Escultura}

\begin{lstlisting}[language=json]
{
    "genero": "escultura",
    "material": "bronce | marmol | madera | hierro",
    "peso_kg": 680,
    "largo_cm": 180,
    "ancho_cm": 98,
    "profundidad_cm": 140
}
\end{lstlisting}

\subsubsection{Fotografía}

\begin{lstlisting}[language=json]
{
    "genero": "fotografia",
    "tecnica": "digital | analogica",
    "papel": "gelatina de plata",       // opcional
    "alto_cm": 30.5,
    "ancho_cm": 40.6,
    "edicion": 3                         // opcional, >= 1
}
\end{lstlisting}

\subsubsection{Cerámica}

\begin{lstlisting}[language=json]
{
    "genero": "ceramica",
    "tipo_arcilla": "porcelana | gres | loza",
    "tecnica_coccion": "bizcocho | esmaltado | raku",
    "peso_kg": 2.5,
    "alto_cm": 30,
    "ancho_cm": 15,
    "profundidad_cm": 15,
    "esmaltado": true                    // opcional
}
\end{lstlisting}

\subsubsection{Orfebrería}

\begin{lstlisting}[language=json]
{
    "genero": "orfebreria",
    "material": "oro | plata | bronce | cobre",
    "tecnica": "filigrana | repujado | fundicion",
    "peso_g": 350,
    "alto_cm": 25,
    "ancho_cm": 12,
    "profundidad_cm": 12,
    "quilates": 18                       // opcional, 0-24
}
\end{lstlisting}

% ===================================================================
\section{Ejemplos completos de petición y respuesta}
% ===================================================================

\subsection{PUT exitoso (Pintura)}

\begin{lstlisting}[language=json, caption=Petición]
PUT /artworks/664a1b2c3d4e5f6a7b8c9d0e
Content-Type: application/json

{
    "genero": "pintura",
    "nombre": "La noche estrellada (actualizada)",
    "artista_id": "artist_001",
    "precio_venta": 18500.00,
    "fecha_creacion": "1889-06-01",
    "foto": "https://ejemplo.com/noche-estrellada.jpg",
    "descripcion": "Oleo sobre lienzo del pintor Vincent van Gogh",
    "tecnica": "oleo",
    "soporte": "lienzo",
    "alto_cm": 73.7,
    "ancho_cm": 92.1
}
\end{lstlisting}

\begin{lstlisting}[language=json, caption=Respuesta 200 OK]
{
    "_id": "664a1b2c3d4e5f6a7b8c9d0e",
    "genero": "pintura",
    "nombre": "La noche estrellada (actualizada)",
    "artista_id": "artist_001",
    "precio_venta": 18500.00,
    "fecha_creacion": "1889-06-01",
    "foto": "https://ejemplo.com/noche-estrellada.jpg",
    "descripcion": "Oleo sobre lienzo del pintor Vincent van Gogh",
    "tecnica": "oleo",
    "soporte": "lienzo",
    "alto_cm": 73.7,
    "ancho_cm": 92.1,
    "status": "disponible",
    "created_at": "2025-01-15T10:30:00",
    "updated_at": "2025-06-01T12:00:00"
}
\end{lstlisting}

\subsection{PUT con ID inexistente (404)}

\begin{lstlisting}[language=json, caption=Petición]
PUT /artworks/aaaaaaaaaaaaaaaaaaaaaaaa
Content-Type: application/json

{ "genero": "pintura", "nombre": "Test", ... }
\end{lstlisting}

\begin{lstlisting}[language=json, caption=Respuesta 404]
{
    "detail": "Obra no encontrada"
}
\end{lstlisting}

\subsection{PUT con body inválido (422)}

\begin{lstlisting}[language=json, caption=Petición]
PUT /artworks/664a1b2c3d4e5f6a7b8c9d0e
Content-Type: application/json

{}
\end{lstlisting}

\begin{lstlisting}[language=json, caption=Respuesta 422]
{
    "detail": [
        { "loc": ["body", "genero"], "msg": "Field required", ... }
    ]
}
\end{lstlisting}

\subsection{DELETE exitoso}

\begin{lstlisting}[language=json, caption=Petición]
DELETE /artworks/664a1b2c3d4e5f6a7b8c9d0e
\end{lstlisting}

\begin{lstlisting}[caption=Respuesta 204 No Content]
( cuerpo vacio )
\end{lstlisting}

\subsection{DELETE con ID inexistente (404)}

\begin{lstlisting}[language=json, caption=Petición]
DELETE /artworks/aaaaaaaaaaaaaaaaaaaaaaaa
\end{lstlisting}

\begin{lstlisting}[language=json, caption=Respuesta 404]
{
    "detail": "Obra no encontrada"
}
\end{lstlisting}

% ===================================================================
\section{Pruebas con curl}
% ===================================================================

A continuación se presentan comandos \texttt{curl} para probar los endpoints.
Reemplazar \texttt{\{ID\}} con un ObjectId real de MongoDB.

\subsection{PUT - Actualizar una pintura}

\begin{lstlisting}[style=bash, caption=Actualizar pintura]
curl -X PUT http://127.0.0.1:8000/artworks/{ID} \
  -H "Content-Type: application/json" \
  -d '{
    "genero": "pintura",
    "nombre": "La noche estrellada",
    "artista_id": "artist_001",
    "precio_venta": 15000.00,
    "fecha_creacion": "1889-06-01",
    "foto": "https://ejemplo.com/noche.jpg",
    "tecnica": "oleo",
    "soporte": "lienzo",
    "alto_cm": 73.7,
    "ancho_cm": 92.1
  }'
\end{lstlisting}

\subsection{PUT - Actualizar una escultura}

\begin{lstlisting}[style=bash, caption=Actualizar escultura]
curl -X PUT http://127.0.0.1:8000/artworks/{ID} \
  -H "Content-Type: application/json" \
  -d '{
    "genero": "escultura",
    "nombre": "El pensador",
    "artista_id": "artist_002",
    "precio_venta": 50000.00,
    "fecha_creacion": "1904-01-01",
    "foto": "https://ejemplo.com/pensador.jpg",
    "material": "bronce",
    "peso_kg": 680,
    "largo_cm": 180,
    "ancho_cm": 98,
    "profundidad_cm": 140
  }'
\end{lstlisting}

\subsection{PUT - Error 404}

\begin{lstlisting}[style=bash, caption=ID inexistente]
curl -X PUT http://127.0.0.1:8000/artworks/aaaaaaaaaaaaaaaaaaaaaaaa \
  -H "Content-Type: application/json" \
  -d '{
    "genero": "pintura",
    "nombre": "Test",
    "artista_id": "x",
    "precio_venta": 100,
    "fecha_creacion": "2024-01-01",
    "foto": "https://ejemplo.com/x.jpg",
    "tecnica": "oleo",
    "soporte": "lienzo",
    "alto_cm": 10,
    "ancho_cm": 10
  }'
\end{lstlisting}

\subsection{DELETE - Eliminar obra}

\begin{lstlisting}[style=bash, caption=Eliminar obra]
curl -X DELETE http://127.0.0.1:8000/artworks/{ID}
\end{lstlisting}

\subsection{DELETE - Error 404}

\begin{lstlisting}[style=bash, caption=ID inexistente]
curl -X DELETE http://127.0.0.1:8000/artworks/aaaaaaaaaaaaaaaaaaaaaaaa
\end{lstlisting}

\subsection{Script automatizado de pruebas}

El proyecto incluye un script PowerShell en \texttt{tests/test\_artworks\_curl.ps1}
que ejecuta una batería completa de pruebas:

\begin{lstlisting}[style=bash, caption=Ejecutar batería de pruebas]
# Sin ID real (prueba solo casos de error 404/422):
.\tests\run_curl_tests.bat

# Con ID real:
.\tests\run_curl_tests.bat 664a1b2c3d4e5f6a7b8c9d0e
\end{lstlisting}

El script ejecuta 11 pruebas:

\begin{enumerate}
    \item PUT Pintura (200)
    \item PUT Escultura (200)
    \item PUT Fotografía (200)
    \item PUT Cerámica (200)
    \item PUT Orfebrería (200)
    \item PUT ID inexistente (404)
    \item PUT ID mal formado (404)
    \item PUT body vacío (422)
    \item DELETE obra existente (204)
    \item DELETE ID inexistente (404)
    \item DELETE ID mal formado (404)
\end{enumerate}

% ===================================================================
\section{Diagrama de flujo}
% ===================================================================

\begin{center}
\begin{minipage}{0.9\textwidth}
\footnotesize
\begin{verbatim}
  Cliente (curl)
       |
       | PUT /artworks/{id}  o  DELETE /artworks/{id}
       v
  +------------------+
  |   FastAPI Router  |
  |  (app/routes/)    |
  +------------------+
       |
       v
  +------------------+
  |    Controller     |
  | (app/controllers/)|
  +------------------+
       |
       | Llamada a MongoDB (Motor)
       v
  +------------------+
  |     MongoDB       |
  |  (Atlas cluster)  |
  +------------------+
       |
       | Respuesta
       v
  +------------------+
  |   Serializacion   |
  |  Pydantic Response|
  +------------------+
       |
       | JSON
       v
  Cliente (curl)
\end{verbatim}
\end{minipage}
\captionof{figure}{Flujo de una petición PUT o DELETE}
\end{center}

% ===================================================================
\section{Buenas prácticas}
% ===================================================================

\begin{itemize}
    \item \textbf{ID válido}: Asegurarse de que \texttt{artwork\_id} sea un
          ObjectId de 24 caracteres hexadecimales. Un ID mal formado produce
          \texttt{404}, no \texttt{400}.
    \item \textbf{Campo \texttt{genero}}: Es obligatorio y debe coincidir
          exactamente con uno de los valores del enum (\texttt{pintura},
          \texttt{escultura}, \texttt{fotografia}, \texttt{ceramica},
          \texttt{orfebreria}).
    \item \textbf{Campos específicos}: Al actualizar, se deben incluir todos
          los campos requeridos del género correspondiente, no solo los que
          se desean modificar (el esquema \texttt{ArtworkCreate} valida el
          objeto completo).
    \item \texttt{updated\_at}: Se asigna automáticamente al momento de la
          actualización. No debe enviarse en el cuerpo.
    \item \texttt{\_id}: Nunca debe incluirse en el cuerpo del PUT. El
          controlador lo elimina por seguridad.
    \item \textbf{PUESTA A PUNTO}: Asegurarse de que el servidor FastAPI
          esté corriendo antes de ejecutar las pruebas curl.
\end{itemize}

% ===================================================================
\section{Conclusión}
% ===================================================================

Los endpoints \texttt{PUT} y \texttt{DELETE} de la API de obras de arte
proporcionan una interfaz sólida y bien tipada para la actualización y
eliminación de registros. La combinación de FastAPI, Pydantic y MongoDB
permite:

\begin{itemize}
    \item Validación rigurosa de datos en el lado del servidor.
    \item Manejo consistente de errores con códigos HTTP estándar.
    \item Discriminación automática del tipo de obra mediante el campo
          \texttt{genero}.
    \item Documentación interactiva generada automáticamente en
          \texttt{/docs} (Swagger UI) y \texttt{/redoc}.
\end{itemize}

\end{document}
